% ============================================================
\section{Introduction}
% ============================================================

\subsection{Cosmic-Ray Neutron Sensing}

The Cosmic-Ray Neutron Sensing (CRNS) method exploits the natural
flux of epithermal neutrons produced when high-energy cosmic-ray
particles interact with the Earth's atmosphere and surface
\citep{zreda2012}.
Secondary neutrons that reach the ground are moderated primarily by
hydrogen-bearing substances; liquid soil water is by far the largest
and most variable hydrogen pool in most natural environments.
A stationary CRNS probe counts epithermal neutrons continuously, and
the count rate $N$ is inversely related to the volumetric soil moisture
$\thetav$ over an area roughly $\sim$\SI{100}{\metre} in radius and
$\sim$\SI{10}{\centi\metre} in depth \citep{desilets2010, kohli2015}.

This footprint is orders of magnitude larger than a conventional
point-scale TDR or FDR sensor yet smaller than remote-sensing products,
making CRNS uniquely suited to bridge the gap between in-situ
hydrology and satellite observations.
The technique is now established globally through networks such as
COSMOS \citep{zreda2012} and TERENO, and is increasingly used for
irrigation scheduling, drought monitoring, and land-surface model
validation.

\subsection{The Need for Pre-Installation Site Characterisation}

A CRNS measurement is not simply a count of neutrons: it is the
result of many simultaneous physical effects that must be corrected
before soil moisture can be retrieved.
Some corrections are \emph{operational} — they must be applied to
each measurement epoch and require ancillary meteorological data
(barometric pressure, atmospheric humidity, incoming cosmic-ray
intensity monitor).
Others are \emph{site-specific constants} that can be determined
once, before installation:

\begin{itemize}[nosep]
  \item \textbf{Topographic correction} $\kT$: departures from a
        flat reference change the soil volume seen by the probe.
        A sensor on a ridge sees less soil than expected; one at the
        bottom of a valley may see more.

  \item \textbf{Muon field-of-view correction} $\kM$: high-energy
        muons (which produce part of the secondary-neutron flux) are
        blocked by surrounding terrain above the horizon, reducing the
        effective production rate.

  \item \textbf{Land-use/land-cover correction} $\kL$: water in
        vegetation, roads, buildings, or standing water within the
        footprint contributes to neutron moderation and must be
        accounted for.

  \item \textbf{Soil hydrogen inventory}: clay-mineral lattice water
        ($\lw$), soil organic carbon (SOC), and bulk density ($\rho_b$)
        set the background hydrogen level against which liquid water
        is detected.
\end{itemize}

Without these corrections, $\Nzero$ — the calibration constant of
the Desilets inversion — is biased, and all subsequent soil moisture
estimates carry a systematic error.
Pre-installation site characterisation is therefore a prerequisite
for accurate, quantitative CRNS measurements.

\subsection{Existing Tools and Gaps}

Several components of CRNS data processing are available as standalone
libraries or web services (e.g.\ the \texttt{crnpy} Python package for
operational corrections, the COSMOS data portal).
However, no openly available, integrated tool exists that:
(i)~downloads all relevant datasets automatically from global open
sources,
(ii)~computes all site-specific correction factors with physically
consistent methods,
(iii)~characterises climate, vegetation, soil, geology, and snow
cover for a given location, and
(iv)~produces a ready-to-use calibration curve and sampling plan.

\subsection{Scope of This Document}

This technical note describes the \emph{CRNS Site Characterization
Pipeline} (\texttt{crns-sitechar}), a Python~3 tool that addresses the
gap above.
Section~\ref{sec:theory} presents the theoretical framework.
Section~\ref{sec:data} summarises all data sources.
Section~\ref{sec:methods} describes the computational methods.
Section~\ref{sec:example} shows a worked example.
Section~\ref{sec:limits} discusses limitations and uncertainties.
Section~\ref{sec:conclusions} concludes.
